% ----------------------------------------------------------
% Metodologia
% ----------------------------------------------------------
\chapter{METODOLOGIA}
\label{cha:metodologia}

% Este capítulo apresenta os procedimentos metodológicos empregados na realização da pesquisa, no desenvolvimento e na avaliação da solução implementada.

% \section{MODALIDADE E ABORDAGEM DA PESQUISA}
% \label{sec:pesquisa}



\section{LEVANTAMENTO DE REQUISITOS} \label{sec:levantamento_requisitos}



O levantamento de requisitos é uma ação fundamental no processo de engenharia de requisitos, ocorrendo principalmente na entre as atividades de comunicação com o cliente e a modelagem do sistema \cite{pressman2011}. 

A etapa de levantamento consiste em identificar o problema, compreender as reais necessidades e expectativas dos diversos \textit{stakeholders} (clientes, usuários finais, gerentes e equipes técnicas) e definir o que o software deve fazer. Ela é fundamental porque estabelece a base entre as necessidades do cliente e a modelagem técnica do sistema, previnindo falhas de comunicação e garantindo que o produto final atenda às expectativas do usuário \cite{pressman2011}.

Neste trabalho, essa etapa teve como objetivo identificar as necessidades e preferências de jogadores e desenvolvedores de jogos digitais, com o intuito de definir os requisitos de interface da plataforma proposta.

Para a coleta de dados, foram elaborados dois questionários on-line, utilizando o Google Forms\footnote{\url{https://docs.google.com/forms/u/0/}}, destinados a públicos distintos: jogadores e desenvolvedores de jogos digitais. A utilização de instrumentos separados permitiu abordar as necessidades específicas de cada perfil de usuário.

Os questionários ficaram disponíveis para resposta durante 14 dias, no período de 22/09 a 06/10. Ao final desse período, foram obtidas 11 respostas de jogadores e 7 respostas de desenvolvedores, totalizando 18 participantes.

Os instrumentos utilizados na coleta de dados, correspondentes aos questionários destinados aos jogadores e aos desenvolvedores, encontram-se disponibilizados na seção de apêndices deste trabalho.


% \subsection{Requisitos funcionais}

% \begin{table}[H]
%     \caption{Requisitos funcionais do perfil desenvolvedor}
% \centering
% \small
% \renewcommand{\arraystretch}{1.4}
% \begin{tabular}{|c|>{\centering\arraybackslash}p{3.2cm}|p{7.5cm}|}
% \hline
% \textbf{ID} & \textbf{Requisito} & \textbf{Descrição} \\
% \hline
% RF01 & Cadastro de jogos & O sistema deve permitir que o desenvolvedor cadastre, altere e remova seus jogos na plataforma. \\
% \hline
% RF02 & Painel centralizado & O sistema deve reunir em um único painel todos os feedbacks recebidos de cada jogo cadastrado. \\
% \hline
% RF03 & Filtro por tipo & O sistema deve permitir que o desenvolvedor filtre os feedbacks por tipo. \\
% \hline
% RF04 & Filtro por urgência & O sistema deve permitir que o desenvolvedor filtre e ordene os feedbacks por nível de urgência. \\
% \hline
% RF05 & Resposta e status & O sistema deve permitir que o desenvolvedor responda publicamente ao feedback e atualize seu status (em análise, implementado ou recusado). \\
% \hline
% RF06 & Visualização em cards & O sistema deve exibir os feedbacks em cards organizados por categoria e urgência. \\
% \hline
% RF07 & Métricas e gráficos & O sistema deve exibir gráficos e métricas sobre os feedbacks. \\
% \hline
% \end{tabular}
% \label{tab:rf-desenvolvedor}
% \legend{Fonte: Autora (2026).}
% \end{table}

% \begin{table}[H]
%     \caption{Requisitos funcionais do perfil jogador}
% \centering
% \small
% \renewcommand{\arraystretch}{1.4}
% \begin{tabular}{|c|>{\centering\arraybackslash}p{3.2cm}|p{7.5cm}|}
% \hline
% \textbf{ID} & \textbf{Requisito} & \textbf{Descrição} \\
% \hline
% RF08 & Envio de feedback & O sistema deve permitir que o jogador envie feedback vinculado a um jogo específico. \\
% \hline
% RF09 & Status do feedback & O sistema deve exibir ao jogador o status do feedback enviado (recebido, em análise, respondido ou implementado). \\
% \hline
% RF10 & Formulário simplificado & O sistema deve oferecer um formulário de envio curto por padrão, com campos opcionais expansíveis. \\
% \hline
% RF11 & Classificação do feedback & O sistema deve permitir que o jogador classifique o feedback por tipo. \\
% \hline
% RF12 & Busca de jogos & O sistema deve permitir que o jogador busque jogos pelo nome. \\
% \hline
% RF13 & Resposta do desenvolvedor & O sistema deve exibir ao jogador a resposta pública do desenvolvedor ao feedback enviado. \\
% \hline
% RF14 & Histórico de feedbacks & O sistema deve permitir que o jogador consulte a lista dos feedbacks que enviou. \\
% \hline
% \end{tabular}
% \label{tab:rf-jogador}
% \legend{Fonte: Autora (2026).}
% \end{table}



% \subsection{Requisitos não funcionais}

% \begin{table}[H]
%     \caption{Requisitos não funcionais comuns aos perfis jogador e desenvolvedor}
% \centering
% \small
% \renewcommand{\arraystretch}{1.4}
% \begin{tabular}{|c|>{\centering\arraybackslash}p{3.2cm}|p{7.5cm}|}
% \hline
% \textbf{ID} & \textbf{Requisito} & \textbf{Descrição} \\
% \hline
% RNF01 & Usabilidade e facilidade de aprendizado &
% A interface deve permitir que os usuários compreendam sua organização, encontrem as funcionalidades e realizem as principais tarefas com facilidade. \\ 
% \hline
% RNF02 & Responsividade & A interface deve adaptar-se a diferentes tamanhos de tela, incluindo dispositivos móveis, tablets e computadores. \\ \hline
% RNF03 & Acessibilidade &
% A interface deve adotar recursos que favoreçam a acessibilidade, como contraste adequado, textos legíveis e elementos de interação identificáveis. \\ 
% \hline
% RNF04 & Consistência visual &
% A interface deve manter padrões consistentes de cores, tipografia, ícones, botões e componentes ao longo das telas. \\ 
% \hline
% RNF05 & Organização da informação & A interface deve apresentar os feedbacks de maneira organizada, com categorias e informações visualmente diferenciadas para facilitar sua identificação. \\ ncionais}

% \begin{table}[H]
%     \caption{Requisitos funcionais do perfil desenvolvedor}
% \centering
% \small
% \renewcommand{\arraystretch}{1.4}
% \begin{tabular}{|c|>{\centering\arraybackslash}p{3.2cm}|p{7.5cm}|}
% \hline
% \textbf{ID} & \textbf{Requisito} & \textbf{Descrição} \\
% \hline
% RF01 & Cadastro de jogos & O sistema deve permitir que o desenvolvedor cadastre, altere e remova seus jogos na plataforma. \\
% \hline
% RF02 & Painel centralizado & O sistema deve reunir em um único painel todos os feedbacks recebidos de cada jogo cadastrado. \\
% \hline
% RF03 & Filtro por tipo & O sistema deve permitir que o desenvolvedor filtre os feedbacks por tipo. \\
% \hline
% RF04 & Filtro por urgência & O sistema deve permitir que o desenvolvedor filtre e ordene os feedbacks por nível de urgência. \\
% \hline
% RF05 & Resposta e status & O sistema deve permitir que o desenvolvedor responda publicamente ao feedback e atualize seu status (em análise, implementado ou recusado). \\
% \hline
% RF06 & Visualização em cards & O sistema deve exibir os feedbacks em cards organizados por categoria e urgência. \\
% \hline
% RF07 & Métricas e gráficos & O sistema deve exibir gráficos e métricas sobre os feedbacks. \\
% \hline
% \end{tabular}
% \label{tab:rf-desenvolvedor}
% \legend{Fonte: Autora (2026).}
% \end{table}

% \begin{table}[H]
%     \caption{Requisitos funcionais do perfil jogador}
% \centering
% \small
% \renewcommand{\arraystretch}{1.4}
% \begin{tabular}{|c|>{\centering\arraybackslash}p{3.2cm}|p{7.5cm}|}
% \hline
% \textbf{ID} & \textbf{Requisito} & \textbf{Descrição} \\
% \hline
% RF08 & Envio de feedback & O sistema deve permitir que o jogador envie feedback vinculado a um jogo específico. \\
% \hline
% RF09 & Status do feedback & O sistema deve exibir ao jogador o status do feedback enviado (recebido, em análise, respondido ou implementado). \\
% \hline
% RF10 & Formulário simplificado & O sistema deve oferecer um formulário de envio curto por padrão, com campos opcionais expansíveis. \\
% \hline
% RF11 & Classificação do feedback & O sistema deve permitir que o jogador classifique o feedback por tipo. \\
% \hline
% RF12 & Busca de jogos & O sistema deve permitir que o jogador busque jogos pelo nome. \\
% \hline
% RF13 & Resposta do desenvolvedor & O sistema deve exibir ao jogador a resposta pública do desenvolvedor ao feedback enviado. \\
% \hline
% RF14 & Histórico de feedbacks & O sistema deve permitir que o jogador consulte a lista dos feedbacks que enviou. \\
% \hline
% \end{tabular}
% \label{tab:rf-jogador}
% \legend{Fonte: Autora (2026).}
% \end{table}



% \subsection{Requisitos não funcionais}

% \begin{table}[H]
%     \caption{Requisitos não funcionais comuns aos perfis jogador e desenvolvedor}
% \centering
% \small
% \renewcommand{\arraystretch}{1.4}
% \begin{tabular}{|c|>{\centering\arraybackslash}p{3.2cm}|p{7.5cm}|}
% \hline
% \textbf{ID} & \textbf{Requisito} & \textbf{Descrição} \\
% \hline
% RNF01 & Usabilidade e facilidade de aprendizado &
% A interface deve permitir que os usuários compreendam sua organização, encontrem as funcionalidades e realizem as principais tarefas com facilidade. \\ 
% \hline
% RNF02 & Responsividade & A interface deve adaptar-se a diferentes tamanhos de tela, incluindo dispositivos móveis, tablets e computadores. \\ \hline
% RNF03 & Acessibilidade &
% A interface deve adotar recursos que favoreçam a acessibilidade, como contraste adequado, textos legíveis e elementos de interação identificáveis. \\ 
% \hline
% RNF04 & Consistência visual &
% A interface deve manter padrões consistentes de cores, tipografia, ícones, botões e componentes ao longo das telas. \\ 
% \hline
% RNF05 & Organização da informação & A interface deve apresentar os feedbacks de maneira organizada, com categorias e informações visualmente diferenciadas para facilitar sua identificação. \\ 
% \hline
% RNF06 & Desempenho & A interface deve apresentar tempos de carregamento e transições entre telas adequados à utilização da plataforma. \\ 
% \hline
% RNF07 & Legibilidade & A interface deve apresentar textos, rótulos, mensagens e informações sobre o status dos feedbacks de forma clara e compreensível. \\ 
% \hline
% RNF08 & Eficiência de interação & A interface deve permitir o envio de feedbacks com poucos passos, oferecendo um formulário conciso e a possibilidade de fornecer informações adicionais. \\ 
% \hline
% \end{tabular}
% \label{tab:rnf-desenvolvedor}
% \legend{Fonte: Autora (2026).}
% \end{table}


% \hline
% RNF06 & Desempenho & A interface deve apresentar tempos de carregamento e transições entre telas adequados à utilização da plataforma. \\ 
% \hline
% RNF07 & Legibilidade & A interface deve apresentar textos, rótulos, mensagens e informações sobre o status dos feedbacks de forma clara e compreensível. \\ 
% \hline
% RNF08 & Eficiência de interação & A interface deve permitir o envio de feedbacks com poucos passos, oferecendo um formulário conciso e a possibilidade de fornecer informações adicionais. \\ 
% \hline
% \end{tabular}
% \label{tab:rnf-desenvolvedor}
% \legend{Fonte: Autora (2026).}
% \end{table}

% \subsection{Diagrama de Caso de Uso UML}




% \section{CARACTERIZAÇÃO DO PROTÓTIPO E ARQUITETURA DO SISTEMA}
% \label{sec:arquitetura}

% \section{TECNOLOGIAS UTILIZADAS}
% \label{sec:tecnologias}

% \section{DESENVOLVIMENTO E IMPLEMENTAÇÃO}
% \label{sec:desenvolvimento}

% \section{TESTES E AVALIAÇÃO DO SISTEMA}
% \label{sec:avaliacao}